\documentclass[11pt]{article}

\usepackage[a4paper,margin=25mm,headheight=15pt]{geometry}
\usepackage[T1]{fontenc}
\usepackage[utf8]{inputenc}
\usepackage{lmodern}
\usepackage{microtype}
\usepackage{amsmath}
\usepackage{booktabs}
\usepackage{array}
\usepackage{graphicx}
\usepackage{xcolor}
\usepackage{fancyhdr}
\usepackage[hidelinks]{hyperref}

\definecolor{Accent}{HTML}{315B7D}
\definecolor{Muted}{HTML}{5B6570}

\hypersetup{pdftitle={Supplementary Results --- Module D}, pdfauthor={Poikela et al.}}

\pagestyle{fancy}\fancyhf{}
\lhead{\small\color{Muted}Supplementary Results --- Module D}
\rhead{\small\color{Muted}Poikela et al.}
\cfoot{\small\thepage}
\renewcommand{\headrulewidth}{0.4pt}
\setlength{\parindent}{0pt}\setlength{\parskip}{0.65em}
\renewcommand{\arraystretch}{1.15}

\newcommand{\Faq}{\textit{F.~aquilonia}}
\newcommand{\Fpo}{\textit{F.~polyctena}}
\newcommand{\Rst}{\ensuremath{R_{\mathrm{st}}}}

\title{\vspace{-3em}\textbf{\color{Accent}Supplementary Results --- Module D}\\[0.3em]
  \Large A structure-corrected two-locus screen for intrinsic incompatibilities}
\author{Poikela et al.}
\date{}

\begin{document}
\maketitle
\vspace{-1.5em}
\begin{center}\small\color{Muted}
This note reports the intrinsic (two-locus) arm. All results are structure-corrected but
pre-simulation: the neutral, recombination-matched null (Methods, \emph{Neutral,
recombination-matched null}) is still required to convert the leads below into
excess-over-neutral statements. Figures and tables are collected here.
\end{center}

% =====================================================================
\section{Most associations reflected technical duplication or broad co-ancestry}\label{srD:scan}

Across $1.38\times10^{6}$ unlinked comparisons in the targeted scan, 276 passed
Benjamini--Hochberg $q<0.01$. Removing 29 likely duplicate or mismapped pairs, consolidating
nearby correlated regions, and excluding associations driven by a single population left
\textbf{110 associations among 83 regions} (18 negative and 92 positive;
Table~\ref{tabRD:cascade}; Fig.~\ref{figRD:network}).

Most of the remaining signal involved low-recombination regions: 66 of the 110 associations
connected two such regions, whereas only 13 connected two regions outside this class. Four
low-recombination regions had at least ten partners. Thus, the scan was dominated by a small
number of broadly associated regions rather than by isolated candidate pairs.

\begin{table}[t]
\centering
\caption{Filtering of the structure-corrected targeted scan. Low recombination denotes a
recombination percentile below 0.10 and is an annotation rather than an exclusion criterion.}
\label{tabRD:cascade}
\small
\begin{tabular}{lrr}
\toprule
Stage & Associations & Regions \\
\midrule
Unlinked focal--partner tests & 1{,}381{,}804 & -- \\
Benjamini--Hochberg $q<0.01$ & 276 & -- \\
After removing likely duplicate mappings & 247 & 156 \\
After regional consolidation & 160 & 127 \\
After leave-one-population-out filtering & 110 & 83 \\
\midrule
\textbf{Final network} & \textbf{110} & \textbf{83} \\
\quad negative / positive & 18 / 92 & -- \\
\quad low--low / mixed / other--other & 66 / 31 / 13 & -- \\
\quad hubs (degree $\geq10$) / low-recombination regions & -- & 4 / 59 \\
\bottomrule
\end{tabular}
\end{table}

\begin{figure}[t]
\centering
\includegraphics[width=0.92\textwidth]{../Figures/moduleD_trans_network.png}
\caption{\textbf{The structure-corrected two-locus network.} Regions are sized by degree;
grey denotes low recombination and gold denotes other regions. Blue edges are positive
(coupling) and orange edges are negative (repulsion). The network contains a large
low-recombination complex, the F33454 co-sorting module, and several isolated candidate pairs.}
\label{figRD:network}
\end{figure}

% =====================================================================
\section{Associations form several multi-chromosomal modules}\label{srD:modules}

Average-linkage clustering separated the network into 23 descriptive co-ancestry modules.
The largest contained 16 regions on 12 chromosomes, followed by modules containing 14, 9,
and 7 regions. Because these modules have distinct genotype patterns, they were examined
separately rather than combined in one ordination. Hierarchical clustering of individuals
within each module showed consistent genotype groupings across its member regions. In the
illustrated module, this pattern extended across six chromosomes
(Fig.~\ref{figRD:modhm}). Overall, 79 of the 110 associations connected regions assigned to
the same module (Fig.~\ref{figRD:modcirc}).

\begin{figure}[t]
\centering
\includegraphics[width=0.94\textwidth]{../Figures/moduleD_module_13_heatmap.png}
\caption{\textbf{A representative multi-chromosomal co-ancestry module.}
Genotype heatmap of a module's member SNPs (9 clusters / 7 meta-nodes across
Chr8/10/11/13/15/21; columns clustered; colour bars: cluster, meta-node, sorting, DI; rows
hierarchically clustered by multilocus genotype). An individual is homozygous-reference
(blue), heterozygous (pale), or homozygous-alternate (red) consistently across all 2{,}926 SNPs
on all six chromosomes. The resulting groups do not coincide exactly with populations.}
\label{figRD:modhm}
\end{figure}

\begin{figure}[t]
\centering
\includegraphics[width=0.86\textwidth]{../Figures/moduleD_module_circos.png}
\caption{\textbf{Within-module associations.} Only edges whose two meta-nodes fall in the same
cross-chromosome correlation module (79 of 110), coloured by module; node ring coloured to
match. Each co-ancestry module reads as its own bundle of links across chromosomes. Outer
track: local LD support ($\mathrm{ld}_w>0.2$) per marker; inner ring: the genome-wide sorting
classification.}
\label{figRD:modcirc}
\end{figure}

% =====================================================================
\section{A small number of exploratory candidate associations remain}\label{srD:candidates}

Thirteen associations connected two regions outside the low-recombination class
(Table~\ref{tabRD:cand}; Fig.~\ref{figRD:candcirc}). Most belonged to an
\emph{F.\ aquilonia} co-sorting module centred on F33454 (Chr10), whose partners occurred on
seven other chromosomes and were predominantly sorted in the same ancestry direction. These
positive associations are more naturally interpreted as coordinated ancestry sorting than as
pair-specific incompatibilities. The strongest negative association joined F11431 and F49480
(Chr3--Chr16; $q=2.2\times10^{-10}$). Neither endpoint was classified as repeatedly sorted,
making this the highest-priority negative-association candidate. Several additional positive
pairs joined regions sorted toward the same parental ancestry. The screen therefore produced
a short exploratory candidate list rather than evidence for a large set of pairwise
incompatibilities.

\begin{table}[t]
\centering
\caption{The 13 associations for which both endpoints lie outside the low-recombination class,
with chromosomes, association sign, false-discovery $q$, and ancestry-sorting class.
The F33454 anchor and its aquilonia-sorted partners form a co-sorting module;
neither endpoint of the strong negative F11431--F49480 association was repeatedly sorted.}
\label{tabRD:cand}
\small
\begin{tabular}{llccl}
\toprule
Pair & Chr & association & $q$ & sorting (a/b) \\
\midrule
F11431--F49480 & 3/16  & negative & $2.2\times10^{-10}$ & unsorted/unsorted \\
F10974--F33454 & 3/10  & negative & $2.7\times10^{-4}$  & aqu/aqu \\
F32802--F33454 & 9/10  & negative & $1.1\times10^{-3}$  & aqu/aqu \\
F11265--F33454 & 3/10  & negative & $6.2\times10^{-3}$  & pol/aqu \\
F15419--F33454 & 4/10  & positive & $2.8\times10^{-11}$ & aqu/aqu \\
F33454--F57178 & 10/19 & positive & $8.8\times10^{-5}$  & aqu/aqu \\
F33454--F41301 & 10/13 & positive & $1.2\times10^{-4}$  & aqu/aqu \\
F33454--F35575 & 10/11 & positive & $2.2\times10^{-4}$  & aqu/aqu \\
F17114--F49429 & 4/15  & positive & $9.2\times10^{-4}$  & pol/pol \\
F15820--F26284 & 4/7   & positive & $2.3\times10^{-3}$  & pol/pol \\
F21099--F33454 & 5/10  & positive & $3.0\times10^{-3}$  & aqu/aqu \\
F22294--F39587 & 6/12  & positive & $4.7\times10^{-3}$  & aqu/aqu \\
F45183--F45504 & 14/14 & positive & $5.0\times10^{-3}$  & unsorted/unsorted \\
\bottomrule
\end{tabular}
\end{table}

\begin{figure}[t]
\centering
\includegraphics[width=0.82\textwidth]{../Figures/moduleD_candidate_circos.png}
\caption{\textbf{Candidate associations on the genome.} Each association is drawn as a band
spanning the two regions it joins (positive blue, negative orange; minimum width for
visibility). The Chr10 anchor F33454 fans out to aquilonia-sorted partners on
Chr3/4/5/9/11/13/19. Neither endpoint of the negative F11431--F49480 association
(Chr3--Chr16) was classified as repeatedly sorted.}
\label{figRD:candcirc}
\end{figure}

% =====================================================================
\section{Interpretive limits}\label{srD:null}

Low recombination does not by itself imply either neutrality or selection. It can preserve
long ancestry blocks, but it can also facilitate the maintenance of coadapted variation.
The empirical scan therefore cannot determine whether the observed modules exceed neutral
expectations. That question requires applying the same procedure to recombination-matched
simulations. Until then, the principal result is that most apparent two-locus associations
can be attributed to likely mapping duplication, broad co-ancestry, or single-population
leverage, leaving a small and transparent exploratory candidate list.

\end{document}
